\documentclass[12pt,a4paper]{article}

% --- Packages ----------------------------------------------------------------
\usepackage[margin=2.5cm]{geometry}
\usepackage{graphicx}
\usepackage{booktabs}
\usepackage{amsmath}
\usepackage{hyperref}
\usepackage{listings}
\usepackage{xcolor}
\usepackage{float}
\usepackage{caption}
\usepackage{parskip}
\usepackage{microtype}
\usepackage{lmodern}

% --- Code listing style -------------------------------------------------------
\lstset{
  basicstyle=\ttfamily\small,
  backgroundcolor=\color{gray!10},
  frame=single,
  breaklines=true,
  columns=fullflexible,
  keepspaces=true,
}

% --- Hyperlink colours --------------------------------------------------------
\hypersetup{
  colorlinks=true,
  linkcolor=blue!60!black,
  urlcolor=blue!60!black,
  citecolor=blue!60!black,
}

% --- Title --------------------------------------------------------------------
\title{%
  \textbf{Project Proposal --- DiT-Based Anomaly Detection\\via Diffusion Reconstruction}\\[6pt]
  \large DATA-MSML 612 --- Deep Learning\\
  University of Maryland
}
\author{Shrikanth Vilvadrinath (121049896) \and Sivram Sahu (121121196) \and Kevin Rathod (121330420) \and Akshay Kamkhalia (121137395)}
\date{March 2026}

% ------------------------------------------------------------------------------
\begin{document}
\maketitle

%===============================================================================
\section{Introduction}
%===============================================================================

Manufacturing defect detection is a critical quality-control challenge.
Human visual inspection is slow, expensive, and inconsistent, yet most supervised
deep learning approaches require large collections of labeled defect images that
rarely exist in practice.

We propose an \textit{unsupervised} approach: train a generative diffusion model
exclusively on defect-free images so that it learns the distribution of normality.
At test time, an image is partially noised and then reconstructed; regions the
model cannot faithfully recover are flagged as anomalous. The key innovation is
our backbone choice: rather than the conventional UNet, we use a
\textbf{Diffusion Transformer (DiT)}~\cite{peebles2023}---making this, to our
knowledge, the first DiT-based anomaly detector evaluated on the MVTec AD
benchmark. The architecture spans both course tracks (Transformer and Diffusion),
and our target is image-level AUROC $>$ 95\%, competitive with leading
diffusion-based detectors~\cite{mousakhan2023}.

%===============================================================================
\section{Background and Related Work}
%===============================================================================

Ho et al.~\cite{ho2020} established the modern diffusion framework: a forward
process adds Gaussian noise over $T$ steps, and a neural network learns to
reverse it via noise prediction. Song et al.~\cite{song2021} showed that the
same model can be sampled deterministically in as few as 50 steps (DDIM), a
$20\times$ speedup over the original 1000-step DDPM schedule.

Reconstruction-based anomaly detection with diffusion was introduced by
Wyatt et al.~\cite{wyatt2022} (AnoDDPM): partially noise a test image, denoise
it, and measure the residual. Mousakhan et al.~\cite{mousakhan2023} pushed
further by adding feature-level comparison alongside pixel residuals, achieving
99.8\% AUROC on MVTec AD. Peebles and Xie~\cite{peebles2023} replaced the UNet
in latent diffusion with a Vision Transformer (DiT), matching or exceeding UNet
quality while offering better scalability. No prior work has combined a DiT
backbone with reconstruction-based anomaly detection---our proposal fills this
gap.

%===============================================================================
\section{Problem Statement}
%===============================================================================

Given only defect-free training images, detect and localize defects in unseen
test images without any labeled anomaly data. This requires solving two
sub-tasks: \textbf{image-level detection} (is the image anomalous?) and
\textbf{pixel-level localization} (which pixels are defective?). The model never
observes anomalous samples during training---a fully unsupervised setting.

%===============================================================================
\section{Dataset and Data Acquisition}
%===============================================================================

We use the \textbf{MVTec Anomaly Detection} dataset~\cite{bergmann2019}, the
standard benchmark for unsupervised industrial anomaly detection. It spans
15 categories (5 textures, 10 objects) with 73 distinct defect types.

\begin{table}[H]
\centering
\caption{MVTec AD dataset summary.}
\begin{tabular}{@{}lcc@{}}
\toprule
\textbf{Split} & \textbf{Images} & \textbf{Ground Truth} \\
\midrule
Train & $\sim$3{,}629 (normal only) & None \\
Test  & $\sim$1{,}725 (mixed)       & Pixel-level masks \\
\midrule
Total & 5{,}354 & 73 defect types across 15 categories \\
\bottomrule
\end{tabular}
\end{table}

The dataset is freely available for academic use (mvtec.com, registration
required, no cost). We train one model per category following the standard
MVTec protocol. Images are resized to $128 \times 128$, normalized to $[-1,1]$,
and augmented with horizontal flips, $\pm10^{\circ}$ rotation, and color jitter.

%===============================================================================
\section{Methodology and Evaluation}
%===============================================================================

\subsubsection*{Architecture}

We train a DiT-S to predict noise under a cosine diffusion schedule, then use
reconstruction error for anomaly scoring:

\begin{lstlisting}
Training (normal images only):
  Forward:   x_t = sqrt(a_t)*x_0 + sqrt(1-a_t)*noise   [cosine schedule]
  Backbone:  DiT-S: Patchify(4x4) -> 12x AdaLN-Attn-MLP blocks -> Unpatchify
  Loss:      MSE(predicted_noise, actual_noise)

Inference (DDIM, 50 steps):
  x_0 -> partial noise (T_partial) -> DDIM reverse -> x_0_hat

Anomaly score (dual-level):
  Pixel:    1 - SSIM(x_0, x_0_hat)
  Feature:  L2(ResNet18(x_0), ResNet18(x_0_hat))  [layers 1-3]
  Final:    alpha * pixel + (1-alpha) * feature
\end{lstlisting}

\subsubsection*{Metrics and Baselines}

All metrics are reported per-category and as the mean across all 15 categories.
We target image-level AUROC $>$ 95\%, pixel-level AUROC $>$ 90\%, and
PRO (Per-Region Overlap, the official MVTec localization metric) $>$ 85\%.
Baselines include PCA reconstruction and a convolutional autoencoder (classical),
plus PatchCore~\cite{roth2022} and Reverse Distillation (SOTA) via the
\texttt{anomalib} library.

\subsubsection*{Ablation Studies}

We plan four ablations to isolate each design choice: (1) DiT vs.\ UNet
backbone, (2) $T_{\text{partial}} \in \{100,200,300,400,500\}$, (3) SSIM vs.\
L2 pixel map, and (4) with vs.\ without feature-level scoring.

%===============================================================================
\section{Framework and Reproducibility}
%===============================================================================

\begin{table}[H]
\centering
\caption{Software and compute stack.}
\begin{tabular}{@{}ll@{}}
\toprule
\textbf{Component} & \textbf{Tool} \\
\midrule
DL Framework       & PyTorch 2.x \\
DiT Reference      & facebookresearch/DiT \\
Diffusion Library  & lucidrains/denoising-diffusion-pytorch \\
Baselines          & anomalib (Intel) \\
Perceptual Metrics & pytorch\_msssim, lpips \\
Compute            & Google Colab Free (T4) + Kaggle Free (P100) \\
\bottomrule
\end{tabular}
\end{table}

All code will be open-source with fixed random seeds (42). Every component in
the stack is free and available for academic use.

%===============================================================================
\begin{thebibliography}{9}

\bibitem{ho2020}
Ho, J., Jain, A., \& Abbeel, P. (2020).
Denoising Diffusion Probabilistic Models.
\textit{NeurIPS}.

\bibitem{peebles2023}
Peebles, W. \& Xie, S. (2023).
Scalable Diffusion Models with Transformers.
\textit{ICCV}.

\bibitem{mousakhan2023}
Mousakhan, A., Brox, T., \& Taber, J. (2023).
Anomaly Detection with Conditioned Denoising Diffusion Models.
\textit{arXiv:2305.15956}.

\bibitem{song2021}
Song, J., Meng, C., \& Ermon, S. (2021).
Denoising Diffusion Implicit Models.
\textit{ICLR}.

\bibitem{bergmann2019}
Bergmann, P., Fauser, M., Sattlegger, D., \& Steger, C. (2019).
MVTec AD -- A Comprehensive Real-World Dataset for Unsupervised Anomaly Detection.
\textit{CVPR}.

\bibitem{wyatt2022}
Wyatt, J., Leach, A., Sherborne, S. M., \& Mayol-Cuevas, W. (2022).
AnoDDPM: Anomaly Detection with Denoising Diffusion Probabilistic Models.
\textit{CVPR Workshop on Visual Anomaly and Novelty Detection}.

\bibitem{roth2022}
Roth, K., Pemula, L., Sohn, J., et al. (2022).
Towards Total Recall in Industrial Anomaly Detection.
\textit{CVPR}.

\end{thebibliography}

\end{document}
